% Figure: induction-cooktop coupling against pan material and drive
% frequency. The eddy-current power deposited in the pan base rises where
% the skin depth is a small fraction of the base thickness and the pan
% resistivity and permeability favour surface dissipation. Ferromagnetic
% steel couples far better than aluminium at the working frequency.
\begin{figure}[htbp]
\centering
\begin{tikzpicture}[>=Latex]
\begin{loglogaxis}[
  width=11.5cm, height=7cm,
  xlabel={drive frequency $f$ ($\si{\hertz}$)},
  ylabel={skin depth $\delta$ ($\si{\milli\meter}$)},
  xmin=1e3, xmax=1e6,
  ymin=0.02, ymax=20,
  legend pos=north east,
  legend cell align=left,
  font=\scriptsize,
  grid=both,
  grid style={enggray!20},
]
% skin depth delta = sqrt(2/(mu sigma omega)) in mm.
% Ferromagnetic steel: mu_r ~ 200, sigma ~ 1e6 S/m. Compute prefactor so
% delta(25 kHz) approx 0.7 mm. delta_mm = C_steel / sqrt(f). At f=25000,
% delta=0.7 => C_steel = 0.7*sqrt(25000) = 110.7.
\addplot[engblue, line width=1.2pt, domain=1e3:1e6, samples=60]
  {110.7/sqrt(x)};
\addlegendentry{ferromagnetic steel pan}
% Aluminium: mu_r=1, sigma ~ 3.5e7 S/m, so delta larger. delta(25 kHz)
% approx 0.53 mm... actually aluminium has larger delta because mu_r=1
% despite higher sigma. Take C_alu so delta(25kHz)=2.1 mm => C=332.
\addplot[engamber, line width=1pt, domain=1e3:1e6, samples=60]
  {332/sqrt(x)};
\addlegendentry{aluminium pan}
% typical 2 mm pan base thickness reference line
\addplot[enggray, dashed, line width=0.8pt, domain=1e3:1e6, samples=2]
  {2};
\node[enggray, anchor=south west, font=\scriptsize] at (axis cs:1.2e3,2.1)
  {$2\,\si{\milli\meter}$ base thickness};
% working band marker at 25 kHz
\addplot[engred, only marks, mark=*, mark size=2pt]
  coordinates {(2.5e4,0.7)};
\node[engred, anchor=west, font=\scriptsize] at (axis cs:2.7e4,0.9)
  {working point $25\,\si{\kilo\hertz}$};
\end{loglogaxis}
\end{tikzpicture}
\caption[Skin depth against frequency for an induction-cooktop pan]{Skin
depth against drive frequency for the base of an induction-cooktop pan, for
a ferromagnetic steel pan and an aluminium pan. The cooktop deposits power
by driving eddy currents in the pan base with an alternating field near
$25\,\si{\kilo\hertz}$; the power concentrates in a surface layer one skin
depth deep. Steel wins twice: its permeability shrinks the skin depth so
the current crowds into a thin high-resistance shell, and that shell has a
higher effective resistance than aluminium's thicker, lower-resistance one,
so more of the drive energy becomes heat in the pan rather than
circulating loss in the coil. At the working frequency the steel skin depth
is a fraction of the base thickness while the aluminium skin depth is a
large fraction of it, which is why a bare aluminium pan couples poorly and
plain cookware carries a ferromagnetic base layer.}
\label{fig:vol04ch08:induction-cooktop}
\end{figure}
